\begin{recipe}{Pasta alla Norma}
\kicker[Hoofdgerecht,Vegetarisch]{Hoofdgerecht · vegetarisch}
\index[register]{Pasta alla Norma}
\index[register]{Aubergine}
\index[register]{Ricotta}
\index[register]{Pasta}

\meta{40 minuten \dvd Siciliaans}

\lettrine{P}{asta} alla Norma komt uit Sicilië en draait om gebakken aubergine,
tomatensaus en ricotta. Neem de tijd voor de aubergine: die mag rustig goudbruin
worden.

\blockrule{Ingrediënten}
\begin{ingredients}
  \ing{400 g}{pasta (penne of rigatoni)}\index[register]{Pasta}
  \ing{1 grote}{aubergine}\index[register]{Aubergine}
  \ingb{olijfolie}
  \ing{1}{knoflookteentje, fijngehakt}
  \ing{150 g}{verse ricotta}\index[register]{Ricotta}
  \ingb{grof zout}
  \ing{1 blik}{tomatenpulp (liefst met basilicumsmaak)}
  \ingb{handvol verse basilicum, grof gesnipperd}
  \ingb{versgemalen peper}
\end{ingredients}

\blockrule{Bereiding}
\tip{Schil de aubergine als kinderen bezwaar hebben tegen de schil.}
\begin{steps}
  \item Was de aubergine en snijd hem ongeschild in blokjes van ± 1 cm. Leg ze op een
        snijplank, strooi er grof zout over en laat 10 minuten staan. Dep ze dan droog
        met keukenpapier.
  \item Bak de aubergineblokjes in olijfolie goudbruin. Doe niet meer blokjes in de pan
        dan er naast elkaar passen en werk in batches. Laat uitlekken op keukenpapier.
  \item Fruit de knoflook op laag vuur in een scheut olijfolie. Zorg dat hij niet bruin
        wordt. Voeg de tomatenpulp toe, breng op smaak met zout en peper en laat
        10 minuten rustig pruttelen.
  \item Kook de pasta al dente. Op een Nederlandse verpakking: kook 1 minuut korter
        dan aangegeven.
  \item Giet de pasta af en meng de tomatensaus, de helft van de ricotta en de basilicum
        erdoor. Voeg de aubergineblokjes toe en roer voorzichtig door.
  \item Schep op en garneer met een lepel ricotta, wat losse basilicumblaadjes en
        versgemalen peper.
\end{steps}
\end{recipe}
